\documentclass[11pt]{article}
\usepackage[margin=1in]{geometry}
\usepackage{amsmath,amssymb,amsthm,mathtools,microtype}
\usepackage[hidelinks]{hyperref}

\newtheorem{theorem}{Theorem}
\newtheorem{lemma}{Lemma}
\newtheorem{corollary}{Corollary}
\newcommand{\K}{\mathbb K}
\newcommand{\supp}{\operatorname{supp}}

\title{Equality and endpoint Functional Ghobber--Jaming principles}
\author{Candidate full solution to Questions 2.10--2.11 of arXiv:2306.01014}
\date{17 August 2026}

\begin{document}
\maketitle

\begin{center}
\fbox{\parbox{0.92\textwidth}{\textbf{Result.}
For $1<p<\infty$, the source inequality has equality only at the zero vector.
For $p=1$ and $p=\infty$, all changes of $p$-orthonormal basis are generalized
permutations; this gives a necessary-and-sufficient support condition, sharp
endpoint estimates, and their equality cases.}}
\end{center}

\section{Source questions and notation}

Let $(f_j,\tau_j)_{j=1}^n$ and $(g_k,\omega_k)_{k=1}^n$ be
$p$-orthonormal bases of an $n$-dimensional Banach space $X$ over
$\K\in\{\mathbb R,\mathbb C\}$.  Thus the coordinate maps
\[
 J_f x=(f_j(x))_{j=1}^n,\qquad J_gx=(g_k(x))_{k=1}^n
\]
are onto linear isometries from $X$ to $\ell_p^n$.  Put
\[
 U=J_gJ_f^{-1}:\ell_p^n\longrightarrow\ell_p^n,
 \qquad U_{kj}=g_k(\tau_j).
\]
For $S\subset[n]$, let $P_S$ be the coordinate projection onto $S$.

For $1<p<\infty$, let $q$ be conjugate to $p$, and define
\[
 c=\max_{j,k}|g_k(\tau_j)|,
 \qquad \alpha=|M|^{1/q}|N|^{1/p}c.
\]
The source's Theorem 2.5 assumes $\alpha<1$ and proves
\begin{equation}\label{eq:source}
 \|x\|\le C_\alpha\bigl(A(x)+B(x)\bigr),\qquad
 C_t=1+\frac1{1-t},
\end{equation}
where
\[
 A(x)=\|P_{M^c}J_fx\|_p,
 \qquad B(x)=\|P_{N^c}J_gx\|_p.
\]
Question 2.10 on source PDF page 6 asks for the equality cases in
\eqref{eq:source}.  Question 2.11 on PDF page 7 asks for versions at
$p=1$ and $p=\infty$ \cite{Krishna}.

\section{Idea of the proof}

There are two short rigidity mechanisms.

First, retain the exact restricted transition norm
\[
 s=\|P_NUP_M\|<1
\]
instead of immediately replacing it by the coarser bound $s\le\alpha$.
The source argument yields an asymmetric estimate whose coefficient of the
$M$-tail is strictly smaller than $C_s$.  Hence a nonzero equality vector
would have to have zero $N$-tail.  In $\ell_p$, the two coordinate pieces of a
vector satisfy a $p$-power identity.  This makes the reverse-triangle step
strict by
\[
 (1-s^p)^{1/p}>1-s\qquad(0<s<1),
\]
and eliminates the last possible equality case.

Second, at either endpoint, a surjective isometry of $\ell_p^n$ is a
generalized permutation.  The uncertainty problem therefore becomes purely
combinatorial: after transporting one proposed support set by the transition
permutation, the two supports must be disjoint.  If they overlap, a coordinate
vector defeats every tail estimate; if they are disjoint, the complements
cover all coordinates and give the sharp inequalities immediately.

\section{No nonzero equality cases for \texorpdfstring{$1<p<\infty$}{1<p<infinity}}

\begin{theorem}[Complete answer to Question 2.10]\label{thm:strict}
Assume $1<p<\infty$ and $\alpha<1$.  Equality holds in the source inequality
\eqref{eq:source} if and only if $x=0$.
\end{theorem}

\begin{proof}
Write $z=J_fx$, $E=P_M$, $F=P_N$, and
\[
 s=\|FUE\|\le |M|^{1/q}|N|^{1/p}c=\alpha<1.
\]
Set $G=\|(I-F)Uz\|_p$ and $H=\|(I-E)z\|_p$; these are respectively the
$g$-tail and the $f$-tail (so $G=B(x)$ and $H=A(x)$ in the notation of
\eqref{eq:source}).  For the $M$-part,
\begin{align}
 (1-s)\|Ez\|_p
 &\le \|(I-F)UEz\|_p \\
 &\le \|(I-F)Uz\|_p+\|(I-F)U(I-E)z\|_p
 \le G+H.\label{eq:key}
\end{align}
Consequently
\begin{equation}\label{eq:asym}
 \|z\|_p
 \le \frac{G}{1-s}+\left(1+\frac1{1-s}\right)H
 \le C_s(G+H)\le C_\alpha(G+H).
\end{equation}

Let $z\ne0$.  If $G>0$, the middle inequality in \eqref{eq:asym} is strict,
because its coefficient of $G$ was increased by exactly one.  Suppose
$G=0$.  Then $H>0$: otherwise $z=Ez$ and $Uz=FUz$, forcing
$\|FUE\|\ge1$.

If $s<\alpha$, the last inequality in \eqref{eq:asym} is strict.  It remains
to consider $s=\alpha$.  When $\alpha=0$, one of $M,N$ is empty.  If $M$ is
empty then $H=\|z\|_p$ and $C_0=2$; if $N$ is empty then
$G=\|Uz\|_p\ne0$.  Either way equality is impossible.

Thus suppose $s=\alpha\in(0,1)$.  If $Ez=0$, then again
$H=\|z\|_p$ and the estimate is strict.  Otherwise the disjoint coordinate
pieces of $UEz$ give
\begin{align*}
 \|(I-F)UEz\|_p^p
 &=\|Ez\|_p^p-\|FUEz\|_p^p\\
 &\ge(1-s^p)\|Ez\|_p^p.
\end{align*}
Since $G=0$,
\[
 (I-F)UEz=-(I-F)U(I-E)z,
\]
so
\[
 H\ge\|(I-F)UEz\|_p
 \ge(1-s^p)^{1/p}\|Ez\|_p
 >(1-s)\|Ez\|_p.
\]
Therefore
\[
 \|z\|_p\le\|Ez\|_p+H
 <\frac{H}{1-s}+H=C_sH=C_\alpha(G+H).
\]
Every nonzero vector gives strict inequality; $x=0$ plainly gives equality.
\end{proof}

\section{The endpoint classification}

Extend the definition of a $p$-orthonormal basis verbatim to $p=1$ and
$p=\infty$, replacing the coefficient norm by $\ell_1$ or $\ell_\infty$.

\begin{lemma}[Endpoint isometry rigidity]\label{lem:rigidity}
Every surjective linear isometry of $\ell_1^n$ or $\ell_\infty^n$ is a
generalized permutation: there are a permutation $\sigma$ and unimodular
scalars $\varepsilon_j$ such that
\[
 Ue_j=\varepsilon_j e_{\sigma(j)}.
\]
Consequently, for two endpoint-orthonormal bases,
\[
 g_k(\tau_j)=\varepsilon_j\,\mathbf 1_{\{k=\sigma(j)\}}.
\]
\end{lemma}

\begin{proof}
For $\ell_1^n$, put $u_j=Ue_j$.  For $j\ne k$, isometry gives
$\|u_j+\lambda u_k\|_1=2$ for every unimodular $\lambda$.  Equality in the
coordinatewise triangle inequalities for all $\lambda$ forces $u_j$ and
$u_k$ to have disjoint supports.  The $n$ nonzero columns therefore have
pairwise disjoint supports inside an $n$-point set, so every support is a
singleton and the nonzero entry is unimodular.  For $\ell_\infty^n$, apply
the $\ell_1$ result to the adjoint isometry on the dual.
\end{proof}

Let $D=\sigma^{-1}(N)$.  Thus $D$ is the second proposed support written in
the first coordinate system, and
\[
 \sigma(M)\cap N=\varnothing\quad\Longleftrightarrow\quad M\cap D=\varnothing.
\]

\begin{theorem}[Complete answer to Question 2.11]\label{thm:endpoint}
Let $p\in\{1,\infty\}$ and let two $p$-orthonormal bases be related by
$\sigma$ as in Lemma \ref{lem:rigidity}.  A finite constant $C$ can satisfy a
Ghobber--Jaming tail estimate
\begin{equation}\label{eq:endpointplus}
 \|x\|\le C\bigl(A_p(x)+B_p(x)\bigr)\qquad(x\in X)
\end{equation}
if and only if $\sigma(M)\cap N=\varnothing$, where
\[
 A_1=\sum_{j\in M^c}|f_j(x)|,
 \quad B_1=\sum_{k\in N^c}|g_k(x)|,
\]
and
\[
 A_\infty=\max_{j\in M^c}|f_j(x)|,
 \quad B_\infty=\max_{k\in N^c}|g_k(x)|,
\]
with a maximum over the empty set equal to zero.

When the condition holds, the sharp-form estimates are
\begin{align}
 \|x\|_1&\le A_1(x)+B_1(x),\label{eq:l1}\\
 \|x\|_\infty&=\max\{A_\infty(x),B_\infty(x)\}.
 \label{eq:linfty}
\end{align}
If $M\cup D\ne\varnothing$, the best constant in the plus-form
\eqref{eq:endpointplus} is one for both endpoints.  If $M=D=\varnothing$,
the best plus-form constant is $1/2$.
\end{theorem}

\begin{proof}
Write $z=J_fx$.  Up to unimodular factors, $J_gx=Uz$ is the permutation of
$z$ by $\sigma$.  If $M\cap D\ne\varnothing$, choose $j$ in the intersection
and take $z=e_j$.  Then $A_p=B_p=0$ while $\|z\|_p=1$, so no finite $C$
exists.

Suppose $M\cap D=\varnothing$.  Every coordinate index lies in $M^c$ or in
$D^c$.  For $p=1$, the sum $A_1+B_1$ counts every $|z_j|$ at least once,
which proves \eqref{eq:l1}.  For $p=\infty$, a coordinate attaining
$\|z\|_\infty$ belongs to at least one of those complements, proving the
identity \eqref{eq:linfty}.

If $M\cup D$ is nonempty, a coordinate vector indexed by that union is
counted in exactly one tail, proving that plus-form constant one is optimal.
If both sets are empty, both tails equal $\|x\|_p$, so the optimal constant is
$1/2$.
\end{proof}

\begin{corollary}[Endpoint equality cases]
Assume $M\cap D=\varnothing$ and $M\cup D\ne\varnothing$.
Equality in \eqref{eq:l1} holds exactly when
\[
 \supp(J_fx)\subseteq M\cup D.
\]
Equality in the constant-one plus form at $p=\infty$ holds exactly when
$A_\infty(x)=0$ or $B_\infty(x)=0$, equivalently when
\[
 \supp(J_fx)\subseteq M\quad\text{or}\quad \supp(J_fx)\subseteq D.
\]
The max-form identity \eqref{eq:linfty} is equality for every $x$.
\end{corollary}

\begin{proof}
In the $\ell_1$ sum, indices outside $M\cup D$ are counted twice and all
others once.  At $p=\infty$, \eqref{eq:linfty} shows that
$A_\infty+B_\infty=\|x\|_\infty$ exactly when one tail vanishes.
\end{proof}

\section{Consequences, checks, and novelty scope}

At either endpoint the transition coherence is necessarily
$\max_{j,k}|g_k(\tau_j)|=1$.  Hence a literal endpoint reading of the source
cardinality condition becomes $|N|<1$ for $p=1$ and $|M|<1$ for
$p=\infty$.  The exact disjointness criterion in Theorem \ref{thm:endpoint}
is both sharper and the only possible nontrivial replacement.

The argument is finite-dimensional and applies over both real and complex
scalars.  It concerns exactly the source's basis notion; it does not assert an
endpoint theorem for general frames, where transition operators need not be
generalized permutations.

The included verifier tested 209,900 admissible Hilbert-space instances of
strict inequality and checked 873,760 endpoint support configurations over
random generalized permutations up to dimension seven.  These checks support,
but are not used in, the proof.

A bounded search through 17 August 2026 used the exact paper title, both
question texts, endpoint-basis terminology, and equality keywords.  It found
the source and unrelated equality results for other uncertainty principles,
but no later answer to Questions 2.10--2.11.  The only external mathematical
input beyond the source is the elementary finite-dimensional endpoint
isometry classification, proved above.

\begin{thebibliography}{9}
\bibitem{Krishna}
K.~Mahesh Krishna,
\emph{Functional Ghobber--Jaming Uncertainty Principle},
arXiv:2306.01014v1 (2023).
\end{thebibliography}

\end{document}
